\section*{Sets and Indices}

\begin{itemize}
    \item $G$: set of goods types, indexed by $g$. For this instance,
    \[
    G=\{A,B,C,D,E\}.
    \]
    \item $N$: ordered set of candidate containers, indexed by $j$. In the implementation,
    \[
    N=\{0,1,\dots,14\},
    \]
    corresponding to a fixed upper bound of 15 candidate containers.
\end{itemize}

\section*{Parameters}

\begin{itemize}
    \item $Q_g$: available quantity of goods type $g$ to be transported.
    \[
    Q_A=120,\quad Q_B=90,\quad Q_C=300,\quad Q_D=90,\quad Q_E=120.
    \]
    \item $w_g$: weight per unit of goods type $g$ (tons/unit).
    \[
    w_A=0.5,\quad w_B=1,\quad w_C=0.4,\quad w_D=0.6,\quad w_E=0.65.
    \]
    \item $C^{\max}$: maximum container capacity (code: \texttt{max\_cap}),
    \[
    C^{\max}=60.
    \]
    \item $C^{\min}$: minimum load for a used container (code: \texttt{min\_load}),
    \[
    C^{\min}=18.
    \]
    \item $D^{\min}$: minimum number of type-$D$ goods required in each used container (code: \texttt{min\_d\_per\_container}),
    \[
    D^{\min}=12.
    \]
    \item $C/A^{\min}$: minimum number of type-$C$ goods required when type-$A$ goods are present (code: \texttt{min\_c\_per\_a}),
    \[
    C/A^{\min}=1.
    \]
    \item $M_A$: big-$M$ constant used to indicate presence of good $A$ in a container. In the code,
    \[
    M_A = Q_A = 120.
    \]
\end{itemize}

\section*{Decision Variables}

\begin{itemize}
    \item $x_{g j}$ (code: \texttt{x[g][j]}): integer number of units of goods type $g$ loaded into container $j$, for all $g\in G$, $j\in N$.
    \[
    x_{g j}\in \mathbb{Z}_{\ge 0}.
    \]
    \item $y_j$ (code: \texttt{y[j]}): binary variable equal to 1 if container $j$ is used, 0 otherwise.
    \[
    y_j\in\{0,1\}.
    \]
    \item $z_j$ (code: \texttt{z[j]}): binary indicator used for the logical condition involving goods $A$ and $C$.
    \[
    z_j\in\{0,1\}.
    \]
\end{itemize}

\section*{Objective}

Minimize the number of used containers:
\[
\min \sum_{j\in N} y_j.
\]

\section*{Constraints}

\begin{align}
\sum_{j\in N} x_{g j} &= Q_g && \forall g\in G
\label{eq:demand}
\\
\sum_{g\in G} w_g x_{g j} &\le C^{\max} y_j && \forall j\in N
\label{eq:maxcap}
\\
\sum_{g\in G} w_g x_{g j} &\ge C^{\min} y_j && \forall j\in N
\label{eq:minload}
\\
x_{D j} &\ge D^{\min} y_j && \forall j\in N
\label{eq:mind}
\\
x_{A j} &\le M_A z_j && \forall j\in N
\label{eq:aindicator}
\\
x_{C j} &\ge C/A^{\min}\, z_j && \forall j\in N
\label{eq:cifa}
\\
y_j &\ge y_{j+1} && \forall j=0,\dots,13
\label{eq:symmetry}
\end{align}

Constraints \eqref{eq:demand} require all available units of each goods type to be transported.  
Constraints \eqref{eq:maxcap}--\eqref{eq:minload} enforce maximum and minimum load on each used container.  
Constraint \eqref{eq:mind} imposes the minimum number of $D$ goods in every used container.  
Constraints \eqref{eq:aindicator}--\eqref{eq:cifa} implement the logic used in the code: if any $A$ goods are loaded in container $j$, then $z_j$ may be set to 1, which forces at least one $C$ good in that container.  
Constraint \eqref{eq:symmetry} is a symmetry-breaking condition that orders used containers before unused ones.

\section*{Notes on Implementation}

\begin{itemize}
    \item The model is implemented as a mixed-integer linear program in \texttt{current\_heuristic.py} and solved with \texttt{GUROBI\_CMD}.
    \item The candidate container set is not endogenous: the code fixes an upper bound of 15 containers and minimizes the number selected through the binary variables $y_j$.
    \item The $A$--$C$ packaging rule is modeled exactly as implemented, not as a stronger equivalence. Specifically, the code uses auxiliary binary variables $z_j$ with
    \[
    x_{A j}\le M_A z_j,\qquad x_{C j}\ge C/A^{\min}\, z_j,
    \]
    so the implication enforced is ``if container $j$ carries any $A$, then it must carry at least one $C$.'' No converse implication is imposed.
    \item There is no explicit linking constraint of the form $x_{g j}\le Q_g y_j$ for all $g,j$. Thus, unused containers are prevented from carrying goods only indirectly through the capacity constraint \eqref{eq:maxcap}, since $C^{\max}y_j=0$ when $y_j=0$.
    \item All shipment quantities are integer, and weights enter linearly as given in the CSV data.
\end{itemize}
